%% ======================================================================
%% 11_declarations.tex -- etica, CRediT, conflicto de intereses, declaracion de IA, agradecimientos, biografias
%% ======================================================================
\section*{Declaraciones éticas}
Este trabajo no involucró participantes humanos, experimentos con animales ni datos recogidos de redes sociales.

\section*{Contribución de los autores (CRediT)}
\textbf{Sebastián Rejas Córdova:} Conceptualización, Metodología, Investigación (integración y verificación), Validación, Redacción -- revisión y edición.
\textbf{George Henry Ramos Astuhuamán:} Software (supervisor en la Raspberry Pi, visión, firmware del octocóptero y de la ESP32-S3, interfaz web), Metodología (modelo cinemático y dinámico, identificación), Análisis formal, Curación de datos, Visualización, Redacción -- borrador original.
\textbf{Johan Nireo Naupay Fabián:} Metodología, Recursos (diseño y auditoría del sistema de alimentación), Investigación (integración), Redacción -- revisión y edición.
\textbf{Hideky Edu Garnique Pisfil:} Investigación (validación en llano), Software (interfaz de usuario), Visualización, Redacción -- revisión y edición.
\textbf{Diego André Miranda Zamora:} Software (control de la base en LabVIEW y enlace UART myRIO--Pi), Investigación (validación en la rampa), Redacción -- revisión y edición.

\section*{Declaración de conflicto de intereses}
Los autores declaran que no tienen intereses financieros en competencia ni relaciones personales conocidas que pudieran haber influido en el trabajo presentado en este artículo.

\paragraph{Financiamiento}
\tocheck{Este trabajo no recibió financiamiento externo; se usaron equipos del laboratorio de UTEC.}

\section*{Declaración sobre el uso de IA generativa y tecnologías asistidas por IA en la redacción}
Durante la preparación de este manuscrito, los autores emplearon de forma puntual las herramientas de inteligencia artificial Claude (Anthropic) y Gemini (Google) únicamente como apoyo en tareas específicas:
\begin{itemize}
  \item revisión y búsqueda preliminar de literatura técnica y antecedentes;
  \item corrección gramatical, ortográfica y de estilo del manuscrito;
  \item depuración puntual de errores en el código desarrollado por el equipo.
\end{itemize}
El diseño mecánico y eléctrico, el modelado matemático, la construcción del prototipo, la programación de los subsistemas, la ejecución de los ensayos y el análisis de resultados fueron realizados por los autores. Después de utilizar estas herramientas, los autores revisaron, editaron y verificaron el contenido frente al código fuente y los registros experimentales, y asumen la plena responsabilidad por el contenido de la publicación.

\section*{Agradecimientos}
Los autores agradecen a los profesores del curso MT5001 Diseño de Sistemas Mecatrónicos (Elvis Jara y Rodrigo Magallanes) y al personal de los laboratorios de la UTEC por el equipamiento (myRIO, kit TETRIX y fuentes de banco) y su orientación.

\section*{Biografías}
\noindent\textbf{Sebastián Rejas Córdova} es estudiante de pregrado de Ingeniería Mecatrónica en la UTEC (Lima, Perú). Sus intereses son el diseño mecánico y la integración de sistemas de robots móviles.

\noindent\textbf{George Henry Ramos Astuhuamán} es estudiante de pregrado de Ingeniería Mecatrónica en la UTEC. Sus intereses son el control embebido, la visión por computadora y la robótica aérea.

\noindent\textbf{Johan Nireo Naupay Fabián} es estudiante de pregrado de Ingeniería Mecatrónica en la UTEC. Sus intereses son la electrónica de potencia y los sistemas de energía para robots.

\noindent\textbf{Hideky Edu Garnique Pisfil} es estudiante de pregrado de Ingeniería Mecatrónica en la UTEC. Sus intereses son las interfaces humano--máquina y la validación experimental.

\noindent\textbf{Diego André Miranda Zamora} es estudiante de pregrado de Ingeniería Mecatrónica en la UTEC. Sus intereses son el control en tiempo real con hardware de NI y la comunicación entre robots.
